\subsection{Overview}

Although Bangladesh has experienced significant economic growth over the past decade, access to business financing remains a major challenge for many Small and Medium Enterprises (SMEs). Entrepreneurs frequently struggle to secure capital for business expansion, equipment purchases, inventory management, and operational growth. While financing opportunities exist through commercial banks and financial institutions, many SMEs remain excluded due to strict eligibility requirements and limited access to ethical investment alternatives.

At the same time, many individuals possess savings that could be invested in productive businesses. However, they often lack access to transparent, trustworthy, and ethically governed investment platforms. This creates a disconnect between businesses seeking capital and investors seeking responsible investment opportunities.

The following subsections describe the primary problems addressed by the proposed system.

%----------------------------------------------------

\subsection{Challenges Faced by SMEs}

Small businesses encounter numerous barriers when attempting to obtain external financing. Traditional financial institutions generally require collateral, lengthy approval procedures, financial history, and extensive documentation before approving loans. New entrepreneurs and micro-businesses often fail to satisfy these requirements despite having commercially viable business ideas.

Furthermore, financing through conventional lending introduces fixed repayment obligations regardless of business performance. This increases financial pressure during periods of low revenue and may discourage entrepreneurs from pursuing business expansion.

Major SME challenges include:

\begin{itemize}
\item Difficulty obtaining financing without collateral.
\item Limited visibility to potential investors.
\item Lengthy loan approval procedures.
\item Dependence on interest-based financing.
\item Limited access to ethical partnership-based investment models.
\end{itemize}

%----------------------------------------------------

\subsection{Challenges Faced by Investors}

Retail investors also encounter significant difficulties when attempting to invest directly in legitimate businesses. Existing investment opportunities frequently require substantial capital, provide limited transparency, or offer insufficient information regarding business performance and associated risks.

Many investors are interested in ethical investment opportunities that comply with Islamic financial principles. However, reliable digital platforms supporting partnership-based investment remain extremely limited.

Major investor challenges include:

\begin{itemize}
\item Lack of verified investment opportunities.
\item Limited transparency regarding business information.
\item Difficulty evaluating investment risk.
\item Fear of fraud and document manipulation.
\item Limited access to Shariah-compliant investment options.
\end{itemize}

%----------------------------------------------------

\subsection{Limitations of Existing Solutions}

Current financing methods only partially address the needs of SMEs and investors.

Commercial banks primarily focus on collateral-based lending rather than partnership investment. Venture capital firms generally invest only in high-growth startups capable of generating significant returns, leaving many small businesses underserved. Conventional crowdfunding platforms often emphasize donations or reward-based campaigns instead of structured investment partnerships.

Moreover, most existing digital platforms provide limited mechanisms for document verification, auditability, and long-term transparency throughout the investment lifecycle.

%----------------------------------------------------

\subsection{Trust and Transparency Issues}

Trust represents one of the most critical barriers preventing direct investment between businesses and individuals.

Investors require confidence that businesses provide authentic documentation, accurate financial information, and transparent reporting. Likewise, businesses require assurance that investors are legitimate participants capable of fulfilling agreed financial commitments.

Without reliable verification mechanisms, secure authentication, immutable audit records, and transparent documentation, both parties remain exposed to fraud, misinformation, and contractual disputes.

The absence of an independent audit mechanism further increases concerns regarding document tampering, unauthorized record modification, and accountability.

%----------------------------------------------------

\subsection{Need for a Software-Based Solution}

The problems identified above demonstrate the need for a secure digital platform capable of connecting verified SMEs and verified investors within a transparent and ethically governed ecosystem.

Rather than replacing financial institutions, the proposed platform focuses on improving trust, transparency, documentation, verification, and investment management through modern Software Engineering practices.

The proposed system incorporates secure authentication, role-based access control, document verification, encrypted storage, cryptographic hashing, and a Permissioned Trust and Audit Layer to strengthen confidence among all stakeholders while maintaining a technically feasible architecture suitable for a university-level Minimum Viable Product (MVP).